\documentclass[11pt]{article}
\usepackage[T1]{fontenc}
\usepackage{lmodern}
\usepackage[margin=1in]{geometry}
\usepackage{amsmath,amssymb,amsthm,booktabs,hyperref,natbib,microtype,enumitem}
\usepackage{tikz}

\theoremstyle{plain}
\newtheorem{theorem}{Theorem}
\newtheorem{proposition}{Proposition}
\newtheorem{lemma}{Lemma}
\newtheorem{corollary}{Corollary}
\theoremstyle{definition}
\newtheorem{definition}{Definition}
\theoremstyle{remark}
\newtheorem{remark}{Remark}

\newcommand{\PASS}{\mathrm{PASS}}
\newcommand{\dhat}{\widehat{\delta}}
\newcommand{\Dcal}{\mathcal{D}}
\newcommand{\Bord}{\succeq_{B}}
\newcommand{\Ccal}{\mathcal{C}}
\newcommand{\cov}{\mathrm{cov}}

\title{Context selection is covering, not ranking:\\
Minimal faithful representations and amortized partial-order selection}

\author{%
  Andriy Bilous \\
  Department of Information Systems and Networks \\
  Lviv Polytechnic National University, Lviv, Ukraine \\
  ORCID: 0009-0006-5467-7932 \\
  \texttt{andriy.bilous@uitware.com}
  \and
  Oleksandr Reminnyi \\
  Department of Computerized Automatic Systems \\
  Lviv Polytechnic National University, Lviv, Ukraine \\
  ORCID: 0009-0005-5119-3695 \\
  \texttt{sasha.reminnyi@gmail.com}
}

\date{}

\begin{document}
\maketitle

\begin{abstract}
A companion paper shows that context sources for a language model are only partially ordered, and
that consequently no single relevance number can rank them faithfully. That is an impossibility
result, and it leaves the constructive question open: if one number is not enough, how many are,
and can they be computed cheaply? We answer both. For sources modeled as deterministic projections
of a latent requirement, we show the Blackwell order is order-isomorphic to the Boolean lattice of
the coordinate sets a source reveals, so a source's \emph{coverage set} is a faithful
representation. Using the order dimension of that lattice, we prove the minimal faithful vector
representation has dimension exactly $d$, the number of latent coordinates; scalar scoring is the
$d=1$ case, which recovers the companion paper's impossibility as a corollary and sharpens it into
an exact requirement. Two structural consequences follow: the non-dominated front is an antichain
and can be exponentially large, so a selection budget is unavoidable, and budgeted selection is
exactly maximum coverage, whose greedy $(1-1/e)$ guarantee is therefore optimal rather than merely
conventional. This reframes context selection as a covering problem over latent requirements rather
than a ranking problem over documents. We give \textsc{Probe-Select}, which realizes the coverage
representation exactly by measurement, and then state the amortization question the reframing makes
precise: coverage vectors are predictable from source text, whereas the scalars that existing
learned retrievers amortize provably are not. [Planned: a non-circular held-out benchmark against
dense, cross-encoder, listwise, and diversity baselines; a learned coverage predictor evaluated by
order preservation and end-task regret; and a naturalistic prevalence study.]
\end{abstract}

% ==========================================================================
\section{Introduction}
\label{sec:intro}
% ==========================================================================

The companion paper~\citep{companion} establishes a negative result: context sources form a
Blackwell partial order, incomparable pairs occur and are verifiable on production models, and
therefore no query-independent scalar score can rank sources faithfully. It also leaves an explicit
gap. Its survey of the design space ends with an empty cell: every method cheap enough for
retrieval scale collapses to a total order, and every method that respects source interaction is
either symmetric redundancy over embedding geometry or as expensive as direct probing. Nothing is
both cheap and faithful to the partial order.

This paper is about that cell. An impossibility result invites exactly one follow-up question, and
the companion paper does not ask it: \emph{if one number is not enough, how many are?} We show the
question has a sharp answer, that the answer changes what kind of object a selection system should
compute, and that the object it identifies is the one thing in this problem that is plausibly
learnable.

\paragraph{From ranking to covering.}
Our starting observation is that for sources modeled as deterministic projections, which is the
companion paper's model, the Blackwell order is not an abstract poset. A source reveals a subset of
the latent coordinates a task requires, and one source dominates another exactly when its
coordinate set contains the other's. The order is therefore order-isomorphic to the Boolean lattice
$(2^{[d]},\subseteq)$, and the natural representation of a source is its \emph{coverage set}. Once
that identification is made, the vocabulary of the problem changes. Selection is not sorting
documents by a score and cutting; it is choosing a set of sources that \emph{covers} the
coordinates the task requires. Context selection is a covering problem, not a ranking problem.

\paragraph{How many numbers are enough.}
The identification is not merely suggestive; it makes the representation question answerable.
Asking for a faithful vector representation is asking for an order-embedding of $2^{[d]}$ into
$\mathbb{R}^m$ under the product order, and the least such $m$ is by definition the order dimension
of the lattice, which is exactly $d$~\citep{dushnik1941,trotter1992}. So $d$ numbers suffice, fewer
than $d$ provably do not, and the scalar pipeline is the $m=1$ special case. The companion paper's
impossibility theorem falls out as the corollary $m=1 < d$, now with a quantitative statement of
what is missing rather than only that something is.

\paragraph{Why the front needs a budget, and why greedy is optimal.}
Two consequences follow immediately and neither is visible from the companion paper's framing.
First, the non-dominated front is an antichain in the lattice, and by Sperner's
theorem~\citep{sperner1928} an antichain can have $\binom{d}{\lfloor d/2 \rfloor}$ elements, so
``keep everything non-dominated'' is not a usable rule at scale and a budget is forced. Second,
under a budget the problem is exactly maximum coverage, so the greedy $(1-1/e)$
guarantee~\citep{nemhauser1978} is not a conventional bound imported for respectability: by
\citet{feige1998} no polynomial-time algorithm improves on it unless $\mathrm{P}=\mathrm{NP}$. The
approximation constant is tight, which is a stronger statement than the companion paper's method
sketch was able to make.

\paragraph{Contributions.}
\begin{enumerate}[leftmargin=1.6em,itemsep=3pt]
  \item \textbf{Coverage representation} (\S\ref{sec:coverage}). The Blackwell order over
    projection sources is order-isomorphic to $(2^{[d]},\subseteq)$; the coverage set is a faithful
    representation (Theorem~\ref{thm:iso}).
  \item \textbf{Exact dimension} (\S\ref{sec:dimension}). The minimal faithful vector
    representation has dimension exactly $d$ (Theorem~\ref{thm:dimension}), with scalar
    impossibility as a corollary (Corollary~\ref{cor:scalar}) and a lower bound on what any learned
    representation must encode.
  \item \textbf{Structure of selection} (\S\ref{sec:structure}). The front is an antichain with a
    Sperner bound (Proposition~\ref{prop:antichain}), and budgeted selection is maximum coverage
    with a tight $(1-1/e)$ constant (Proposition~\ref{prop:maxcover}).
  \item \textbf{Exact realization} (\S\ref{sec:probeselect}). \textsc{Probe-Select} computes the
    coverage order by measurement, with monotone safety and strict generalization of top-$k$.
  \item \textbf{The amortization question, made precise} (\S\ref{sec:amortize}). The dimension
    bound explains why existing learned retrievers cannot reach the empty cell (they amortize a
    scalar) and identifies coverage-vector prediction as the target that can.
\end{enumerate}

\paragraph{Relation to the companion paper.}
We import and do not re-derive: the experiment model, the fixed model rule $\delta_M$, the
decision-restricted surrogate $\dhat_{\Dcal}$, the finite-sample verification machinery, and the
verified anti-monotonicity phenomenon are all established in~\citet{companion} and used here as
given. This paper contributes the representation theory over that model, the structural
consequences for selection, and the amortization program. Where the companion paper asks whether
scalar selection is broken, this paper asks what to compute instead.

% ==========================================================================
\section{Preliminaries}
\label{sec:prelim}
% ==========================================================================

We restate only what the theorems need; see~\citet{companion} for the full development.

A task $T$ carries a verifier $v_T$, and $\PASS_M(W,T)$ is the pass rate of a fixed model rule
$\delta_M$ given context source $W$. The latent requirement is a vector
$S=(S_1,\dots,S_d)$ of \emph{coordinates}: the private facts a solution must respect. Coordinates
are non-degenerate and mutually independent under the prior. A \emph{projection source} $W_A$, for
$A \subseteq [d]$, reveals $(S_i)_{i \in A}$ and nothing else; write $\cov(W_A) = A$ for its
coverage set. We write $W \Bord W'$ for the Blackwell (garbling) order and take the
Blackwell--Sherman--Stein equivalence as given~\citep{blackwell1953,torgersen1991}.

Verification of empirical dominance and incomparability from finitely many measurements, via exact
Clopper--Pearson intervals with a union bound at level $\eta$, is the companion paper's machinery
and is used unchanged in \S\ref{sec:probeselect}.

% ==========================================================================
\section{The coverage representation}
\label{sec:coverage}
% ==========================================================================

\begin{theorem}[Coverage isomorphism]\label{thm:iso}
For projection sources, $W_A \Bord W_B$ if and only if $B \subseteq A$. Hence $A \mapsto W_A$ is an
order isomorphism from $(2^{[d]},\subseteq)$ onto the Blackwell order restricted to projection
sources.
\end{theorem}

\begin{proof}
If $B \subseteq A$, the map discarding the coordinates in $A \setminus B$ is a deterministic Markov
kernel carrying $W_A$'s signal to $W_B$'s, so $W_A \Bord W_B$. Conversely, suppose $B \not\subseteq
A$ and pick $i \in B \setminus A$. A garbling of $W_A$ is measurable with respect to
$\sigma((S_j)_{j \in A})$, which by independence and non-degeneracy of the coordinates is
independent of $S_i$; no such kernel reproduces $W_B$'s signal, which determines $S_i$. Hence
$W_A \not\Bord W_B$. Antisymmetry and transitivity transfer from set inclusion.
\end{proof}

The content of Theorem~\ref{thm:iso} is that the order carries no structure beyond containment of
coverage sets. Everything below is a consequence of working in a Boolean lattice rather than in an
arbitrary poset.

\subsection{How many numbers a faithful representation needs}
\label{sec:dimension}

\begin{definition}[Faithful representation]\label{def:faithful}
A map $\varphi$ from sources to $\mathbb{R}^m$ is \emph{faithful} if for all sources $W,W'$,
\[
  W \Bord W' \iff \varphi(W) \ge \varphi(W') \ \text{componentwise.}
\]
\end{definition}

Faithfulness is the property a selection system needs: it may compare representations instead of
sources without ever concluding a dominance that does not hold, or missing one that does.

\begin{theorem}[Minimal faithful dimension]\label{thm:dimension}
The least $m$ admitting a faithful representation of the projection sources on $d$ coordinates is
exactly $m = d$. The coverage indicator $\varphi(W_A) = \mathbf{1}_A \in \{0,1\}^d$ attains it.
\end{theorem}

\begin{proof}
\emph{Attainment.} $\mathbf{1}_B \le \mathbf{1}_A$ componentwise iff $B \subseteq A$, which by
Theorem~\ref{thm:iso} holds iff $W_A \Bord W_B$; so $\varphi = \mathbf{1}_{(\cdot)}$ is faithful
with $m=d$.

\emph{Lower bound.} A faithful $\varphi$ is precisely an order-embedding of $(2^{[d]},\subseteq)$
into $\mathbb{R}^m$ with the product order. By Ore's characterization, the least $m$ for which a
finite poset embeds into a product of $m$ chains is its order dimension~\citep{trotter1992}, and
the order dimension of the Boolean lattice $2^{[d]}$ is $d$~\citep{dushnik1941}. Hence $m \ge d$.
\end{proof}

\begin{corollary}[Scalar impossibility, quantified]\label{cor:scalar}
A scalar score is faithful if and only if $d \le 1$. For $d \ge 2$ every scalar mis-orders some
pair of sources, recovering the companion paper's incomparability theorem; moreover no
representation of dimension $m < d$ repairs it.
\end{corollary}

\begin{remark}[What this adds to the companion paper]
The companion paper proves that \emph{some} faithful scalar fails to exist. Theorem~\ref{thm:dimension}
says how far from sufficient a scalar is, and the gap is not a constant: it grows with the number
of distinct latent requirements in play. This converts an impossibility into a design specification,
namely that a selection system must carry $d$ numbers per source, and it is what makes the
amortization question of \S\ref{sec:amortize} well posed rather than merely hopeful.
\end{remark}

\subsection{Consequences for selection}
\label{sec:structure}

\begin{proposition}[The front is an antichain, and can be large]\label{prop:antichain}
The set of non-dominated sources in a candidate pool is an antichain in $(2^{[d]},\subseteq)$. In
the worst case its size is $\binom{d}{\lfloor d/2 \rfloor}$, attained by the middle
layer~\citep{sperner1928}.
\end{proposition}

\begin{proof}
Non-domination is mutual incomparability, which by Theorem~\ref{thm:iso} is mutual non-containment,
the definition of an antichain. Sperner's theorem gives the extremal size.
\end{proof}

So ``return the non-dominated set'' is well defined but not by itself an algorithm: the front is
exponentially large in the worst case. A budget is forced, and with it the question of which
members of the front to keep, which the next proposition settles.

\begin{proposition}[Budgeted selection is maximum coverage, and greedy is optimal]\label{prop:maxcover}
Selecting $k$ sources from a front to maximize the number of distinct covered coordinates is an
instance of maximum coverage. The greedy rule attains $(1-1/e)$ of the
optimum~\citep{nemhauser1978}, and no polynomial-time algorithm attains $(1-1/e)+\epsilon$ for any
$\epsilon>0$ unless $\mathrm{P} = \mathrm{NP}$~\citep{feige1998}.
\end{proposition}

\begin{remark}[Why the constant matters here]
A companion-paper draft of this method presented $(1-1/e)$ as a standard guarantee and claimed no
novelty in it, which understated the situation. Under Theorem~\ref{thm:iso} the objective
\emph{is} coverage, not a submodular surrogate chosen for convenience, and therefore
\citet{feige1998} applies: the greedy step is not merely adequate, it is the best any efficient
selector can do. Diversity heuristics defined over embedding
geometry~\citep{carbonell1998,ye2023ceil} optimize a different objective and inherit no such
optimality.
\end{remark}

% ==========================================================================
\section{\textsc{Probe-Select}: realizing the coverage order by measurement}
\label{sec:probeselect}
% ==========================================================================

Theorem~\ref{thm:dimension} says what to compute. This section gives the estimator that computes it
exactly, at probe cost; \S\ref{sec:amortize} asks whether it can be predicted instead.

\paragraph{Verified empirical dominance.}
Fix $\delta_M$ and a probe battery $\Dcal$. Write $W_i \trianglerighteq W_j$ when the companion
paper's verification jointly establishes, at confidence $\ge 1-\eta$, that
$\dhat_{\Dcal}(W_i,W_j)=0$ up to a tolerance $\alpha_{\mathrm{tol}}$ and that
$L_{\Dcal}(W_j,W_i)>0$. Mutually non-dominated sources are verified incomparable and constitute the
empirical front. Under Theorem~\ref{thm:iso}, $\trianglerighteq$ is the measured image of
coordinate containment, and the probe battery plays the role of the coordinate basis.

\noindent\textbf{Algorithm 1 (\textsc{Probe-Select}).} \emph{Input:} candidates $\Ccal$, probe
battery $\Dcal$, budget $k$, model $\delta_M$, level $\eta$.
\begin{enumerate}[leftmargin=2.2em,itemsep=2pt]
  \item Measure $\PASS(W,T)$ for all $W \in \Ccal$, $T \in \Dcal$ ($n$ draws per cell); form the
    Clopper--Pearson intervals once.
  \item Build $\trianglerighteq$ from those intervals; no extra model calls.
  \item Prune every $W$ verified-dominated by some $W' \in \Ccal$.
  \item Let $P$ be the surviving front.
  \item If $|P| \le k$ return $P$; else greedily add the survivor with the largest verified
    marginal PASS gain on probe tasks not yet covered, breaking ties by smallest $\dhat_{\Dcal}$ to
    the chosen set, until $k$ are selected.
\end{enumerate}

\begin{proposition}[Monotone safety]\label{prop:safe}
\textsc{Probe-Select} never returns a source verified-dominated by an available source. With
probability $\ge 1-\eta$ it never discards a source that strictly wins a probe task in favour of
one that is nowhere better than it by more than $\alpha_{\mathrm{tol}}$.
\end{proposition}

\begin{proof}
Step 3 removes only verified-dominated sources and step 5 selects from the survivors. On the joint
$(1-\eta)$ event a pruned $W$ has no probe task on which it beats its dominator by more than
$\alpha_{\mathrm{tol}}$, and the dominator is retained or Pareto-equivalent to a retained source.
\end{proof}

\begin{proposition}[Strict generalization of top-$k$]\label{prop:generalize}
If $\trianglerighteq$ totally orders $\Ccal$, the front is a chain and \textsc{Probe-Select}
returns its top $k$ elements, recovering scalar top-$k$. It therefore extends the scalar pipeline
and differs from it only on incomparable pairs, which by Corollary~\ref{cor:scalar} are exactly the
pairs a scalar cannot handle.
\end{proposition}

\paragraph{Cost.}
Building $\trianglerighteq$ on $m=|\Ccal|$ candidates costs $m \cdot |\Dcal| \cdot n$
verifier-gated calls (the pairwise comparisons reuse one per-(source,\,task) PASS table, so cost is
linear in $m$) plus $O(m^2)$ interval comparisons. This is curation scale, not corpus scale, and it
is the cost the next section attacks.

% ==========================================================================
\section{Amortization: the cheap-and-faithful cell}
\label{sec:amortize}
% ==========================================================================

\paragraph{Why prior amortization cannot reach it.}
Learned retrievers already amortize ``value to the model'': REPLUG and REALM train a retriever on
the language model's own signal, paying an offline training cost to buy a cheap query-time
score~\citep{shi2023replug,guu2020realm}. By Corollary~\ref{cor:scalar} this cannot be faithful for
$d \ge 2$, and the reason is now precise rather than rhetorical: they amortize an object of
dimension $1$, and the target has dimension $d$. Multi-vector late
interaction~\citep{khattab2020colbert} carries many numbers per document but aggregates to a single
query-document score before selection, so the selection-time object is again one-dimensional. The
failure is not insufficient training; it is a dimension deficit.

\paragraph{The target this identifies.}
Theorem~\ref{thm:dimension} names the object to learn: a map $\widehat{\varphi}$ from source text
to a coverage vector in $\mathbb{R}^d$, whose induced product order approximates
$\trianglerighteq$. This is the amortization the design space is missing, and unlike the scalar case
it is not excluded by any result we know. Training labels are exactly what Algorithm~1 already
produces, so the probe tables computed for exact selection are the supervision for the predictor.

\paragraph{What would count as success.}
Three measurements, in increasing strength: per-coordinate recovery of $\widehat{\varphi}$ against
verified coverage; \emph{order preservation}, the fraction of candidate pairs whose dominance or
incomparability verdict is preserved by $\widehat{\varphi}$; and \emph{end-task regret}, the PASS
gap between selecting with $\widehat{\varphi}$ alone and selecting with full probing. Order
preservation is the decisive one, since it is the property Definition~\ref{def:faithful} requires
and the property a scalar provably cannot have.

\begin{remark}[Honest statement of the bet]
Whether coverage vectors are predictable from source text at useful accuracy is an empirical
question and we do not claim it in advance. The theory establishes that the target is the right one
and that no lower-dimensional target can work; it does not establish learnability. A negative
outcome is informative and reportable: it would locate the obstruction in prediction rather than in
representation, which is itself the answer to the companion paper's open problem.
\end{remark}

% ==========================================================================
\section{Experimental program}
\label{sec:experiments}
% ==========================================================================

[PLANNED. This section is the empirical content of the paper and is not yet run.]

\begin{enumerate}[leftmargin=2em,itemsep=3pt]
  \item \textbf{Non-circular selection benchmark (curation scale).} A pool of $8$--$15$ candidate
    sources with held-out evaluation tasks \emph{disjoint} from the probe battery; endpoint is
    end-task PASS of the assembled context. Baselines: dense bi-encoder top-$k$, a cross-encoder
    reranker, an LLM listwise reranker~\citep{sun2023rankgpt}, MMR~\citep{carbonell1998}, and the
    superset heuristic, all at matched probe budget. Disjointness is the load-bearing design
    constraint: overlap between probe and evaluation tasks would make the comparison circular.
  \item \textbf{Coverage-dimension estimate.} Recover the effective $d$ of a candidate pool from
    the verified order and compare the observed front size against the Sperner bound of
    Proposition~\ref{prop:antichain}.
  \item \textbf{Learned coverage predictor.} Train $\widehat{\varphi}$ on probe tables; report
    per-coordinate recovery, order preservation, and end-task regret (\S\ref{sec:amortize}).
  \item \textbf{Partial-order re-ranking head at retrieval scale.} Retriever shortlist $\to$
    \textsc{Probe-Select} on the top-$m$, against reranker-only pipelines on a naturalistic corpus,
    reporting the \emph{prevalence} of verified incomparability among top-$m$ candidates. This
    answers the companion paper's open prevalence question.
  \item \textbf{Sensitivity.} Probe-battery size $|\Dcal|$, per-cell $n$, and $\eta$ ablations;
    the cost--quality frontier.
\end{enumerate}

% ==========================================================================
\section{Scope and limitations}
\label{sec:limitations}
% ==========================================================================

The representation results hold in the deterministic-projection model, which is the companion
paper's model and the setting in which its evidence was collected. For general stochastic sources
the Blackwell order is garbling rather than containment, and no finite-dimensional product-order
embedding is faithful in general; Theorems~\ref{thm:iso} and~\ref{thm:dimension} should be read as
sharp statements about projection sources, not about arbitrary experiments. The coordinate count
$d$ is a property of a task family together with a candidate pool, not a universal constant, and
estimating it is itself an experimental question (item 2 above). Finally, $\trianglerighteq$ is
verified against a probe battery, so it inherits the companion paper's decision-restricted reading:
conclusions are relative to the fixed model rule $\delta_M$ and to the coordinates the battery can
resolve.

% ==========================================================================
\section*{Reproducibility}
% ==========================================================================
The verification machinery, harness, and all measured cells referenced here are released with the
companion paper at \url{https://github.com/abilous-ti/blackwell-llm-context}.

\bibliographystyle{plainnat}
\bibliography{blackwell-paper}

\end{document}
